% 核心观点：反馈只有把新增的任务相关证据转化为及时、有效的动作改变，才改善任务完成概率。
% 叙述逻辑：执行变化提出新增证据的需求 -> D1刻画信息的潜在价值 -> D2调用C的信息保留/利用分解 ->
% D3分析实际采用，并调用B1性能差分连接闭环任务结果；不是为历史与残差另建独立分析。
% C到D：关系界须对反馈联合表示及执行时刻的关系成立，不沿用基础计划时刻的几何误差。
% 边界：不把GRU/MLP误差差值当成理想历史价值，不把残差幅度或动作MSE改善当成任务改进。
% \subsection{历史、修正与任务收益}
\subsection{History, Corrections, and Task Benefit}
\label{sec:feedbackanalysis}
% 新增执行证据的收益取决于原先可用的信息，以及所生成的修正是否被实际采用。
The benefit of additional execution evidence depends on the information already available and on whether the resulting correction is actually adopted.

% \subsubsection{新观测与历史的决策价值}
\subsubsection{Decision Value of New Observations and History}
% 以下以动作块执行为具体场景，比较已有计划给出的基础动作与利用新增执行信息生成的修正动作。
We take action-chunk execution as the concrete setting, comparing a base action supplied by an existing plan with a corrected action generated using additional execution information.
% 在固定评价分布 $\mu$ 上，以 $t$ 表示目标执行时刻，将该控制步的动作位置称为执行槽位。
Under a fixed evaluation distribution $\mu$, let $t$ denote the target execution time, and call the action position at that control step an execution slot.
% 基础策略为该槽位选中的动作记为
Denote the action selected by the base policy for this slot by
$B=\pi_0(x)$.
% 这里 $\pi_0$ 为包含计划更新、槽位选择及预先规定回退流程的完整基础执行规则。
Here, $\pi_0$ is the complete base execution rule, including plan updates, slot selection, and predefined fallback procedures.
% 本节参考策略保留相同的基础与反馈推理、接收和调度流程，仅在执行输出端屏蔽残差。
The reference policy retains the same base and feedback inference, reception, and scheduling procedures, but masks the residual at the execution output.
% 计划信息 $\mathcal F_p$ 包含生成和选择该动作所使用的观察与计划上下文，
The planning information $\mathcal F_p$ contains the observations and plan context used to generate and select this action,
% 以及触发回退所需的运行记录，使 $B$ 对 $\mathcal F_p$ 可测。
together with runtime records needed to trigger fallback, so that $B$ is $\mathcal F_p$-measurable.
% 当前信息 $\mathcal F_c$ 在保留计划信息的基础上加入时刻 $t$ 最新可用的观测和本体信息；
The current information $\mathcal F_c$ retains the planning information and adds the latest observations and proprioception available at time $t$.
% 完整因果历史 $\mathcal F_h$ 进一步包含此前的执行记录，
The complete causal history $\mathcal F_h$ additionally contains previous execution records,
% 满足 $\mathcal F_p\subseteq\mathcal F_c\subseteq\mathcal F_h$。
with $\mathcal F_p\subseteq\mathcal F_c\subseteq\mathcal F_h$.
% 运行时信息仅包含已获得的记录及时间戳。
Runtime information includes only records and timestamps already available.

% 新观测和额外历史分别使最优平均失败概率降低的幅度为
The reductions in optimal mean failure probability due to new observations and additional history are, respectively,
\begin{align}
 \Delta_{\mathrm{new}}
 &=\mathcal L_\mu^*(\mathcal F_p)-\mathcal L_\mu^*(\mathcal F_c),\nonumber\\
 \Delta_{\mathrm{hist}}
 &=\mathcal L_\mu^*(\mathcal F_c)-\mathcal L_\mu^*(\mathcal F_h).
 \label{eq:taskinformationgain}
\end{align}
% 信息嵌套使两项均非负，因为较丰富的信息允许复用原动作规则。
Both terms are nonnegative by information inclusion: a rule using less information remains available when more information is provided.
% 一个严格正历史价值的充分情形是：相同当前观测和计划下，两种等概率模式要求相反动作，历史识别模式，正确与错误动作的失败概率差为c>0。
% 可选动作仅为两个相反动作；相同当前信息指全部F_c，包括本体、计划和运行记录；评价分布仅在末次决策上。
For a final-decision distribution $\mu$ in a two-action setting, consider two equally likely modes with identical $\mathcal F_c$, including observations, proprioception, plan context, and runtime records.
If history distinguishes the modes and they require opposite actions, with an incorrect choice increasing failure probability by $0<c\le1$, then $\Delta_{\mathrm{hist}}=c/2$ (Appendix~\ref{app:historyvalue}).
% 该例不是声称所有历史都有价值；若当前信息已决定最优动作，历史价值可为零。
If current information already determines an optimal action, the history value can instead be zero.
% 有限动作--响应历史的编码可能丢失决策所需信息，该损失由后续表示损失项刻画。
Encoding a finite action--response history may discard decision-relevant information; the resulting loss is accounted for by the representation term below.

% \subsubsection{从信息价值到实际修正}
\subsubsection{From Information Value to Implemented Corrections}
% 新信息的潜在价值不等于候选修正的实际收益；沿用C的分解，还须扣除反馈表示的信息损失与决策实现误差。
The potential value of new information is not yet the benefit of a candidate correction.
Applying the decomposition in~\eqref{eq:taskrepresentationdecomposition} to the feedback representation also accounts for information loss and decision-realization error.
% 实际反馈表示 $Z_t^{\mathrm{fb}}$ 编码基础计划、缓存查询、RGB、因果历史及候选选择所需运行记录，
The implemented feedback representation $Z_t^{\mathrm{fb}}$ encodes the base plan, cached queries, RGB inputs, causal history, and runtime records required for candidate selection,
% 生成信息 $\mathcal U=\sigma(Z_t^{\mathrm{fb}})\subseteq\mathcal F_h$。
generating $\mathcal U=\sigma(Z_t^{\mathrm{fb}})\subseteq\mathcal F_h$.
% 在基础计划有效时，基础查询对应计划观测时刻 $\tau_p$，并在计划内保持不变；
When the base plan is valid, its queries correspond to the planning observation time $\tau_p$ and remain unchanged within the plan.
% 候选修正使用最新观测时刻 $\tau_o$ 的图像，作用于槽位 $t$，其中 $\tau_p\le\tau_o\le t$。
Candidate corrections use images from the latest observation time $\tau_o$ and target slot $t$, where $\tau_p\le\tau_o\le t$.
% 编码过程中的信息丢失与观测延迟共同限制该表示所保留的信息。
Information loss during encoding and observation delay jointly limit the information retained by this representation.

% 由该表示生成候选指令 $\widehat u=B+r\in\mathcal A$，
The representation produces a candidate command $\widehat u=B+r\in\mathcal A$,
% 其中 $r$ 为包含限幅等后处理的目标槽位残差。
where $r$ is the residual for the target slot after postprocessing, including magnitude bounding.
% 基础动作的决策实现差距为
The decision-realization gap of the base action is
$\epsilon_{\mathrm{base}}=\mathcal L_\mu(B)-\mathcal L_\mu^*(\mathcal F_p)$.

% \begin{proposition}[候选反馈的任务收益分解]
\begin{proposition}[Decomposition of candidate feedback benefit]
\label{prop:taskfeedback}
% 在同一 $\mu$、$\pi_0$ 和动作集合下，将采用候选修正后平均失败概率的下降称为候选反馈收益，满足
Under the same $\mu$, $\pi_0$, and action set, define the candidate feedback benefit as the reduction in mean failure probability from applying the candidate correction. It satisfies
\begin{align}
 \Gamma_\mu
 &:=\mathcal L_\mu(B)-\mathcal L_\mu(\widehat u)\nonumber\\
 &=\epsilon_{\mathrm{base}}+\Delta_{\mathrm{new}}+\Delta_{\mathrm{hist}}
   -\Delta_{\mathrm{rep}}(\mathcal U)-\epsilon_{\mathrm{dec}}(\widehat u).
 \label{eq:taskfeedbackdecomposition}
\end{align}
\end{proposition}
% 证明见附录~\ref{app:feedbackproofs}。
The proof is given in Appendix~\ref{app:feedbackproofs}.
% 决策实现差距包含残差幅度/容量限制、学习误差及监督目标差异。
The decision-realization gap includes residual magnitude and capacity limits, learning error, and the mismatch between action supervision and task consequences.
% 同分布同槽位的残差动作误差恒等式仅诊断修正方向与幅度，不能替代这里的任务收益。
The matched-distribution, matched-slot action-error identity~\eqref{eq:residualgain} diagnoses correction direction and magnitude, but does not replace this task-benefit criterion.

% \subsubsection{时序对齐与闭环任务收益}
\subsubsection{Temporal Alignment and Closed-Loop Task Benefit}
% 候选的决策优势只有通过实际采用才影响轨迹；以下结合采用规则与B1性能差分。
A candidate's decision advantage affects the trajectory only through its actual adoption.
We connect the adoption rule to task performance using the performance-difference identity~\eqref{eq:pd}.
% r是执行机制实际选中并限幅的目标槽位候选，无候选取0；候选必须及时到达、未过期且匹配有效计划与槽位。
Let $r$ be the bounded target-slot candidate actually selected by the deterministic execution rule, or zero if none is available.
A correction must arrive in time, remain unexpired, and match the valid plan and execution slot.
% 修正采用指示量 $\Lambda\in\{0,1\}$ 在满足这些条件时取 $1$，否则取 $0$。
The correction-adoption indicator $\Lambda\in\{0,1\}$ equals $1$ when these conditions hold and $0$ otherwise.
% 无有效修正时执行 $B$；基础计划也不可用时，$B$ 为同一历史下预先规定的回退指令，
Without a valid correction, $B$ is executed. If the base plan is also unavailable, $B$ is the predefined fallback command for the same history,
% 此时置 $\Lambda=0$，并将候选残差约定为 $r=0$。
with $\Lambda=0$ and the candidate residual set to $r=0$.
% 实际执行指令为 $u_{\mathrm{exec}}=B+\Lambda r$，
The executed command is $u_{\mathrm{exec}}=B+\Lambda r$.
% 在执行情形 $x$ 下，候选修正相对基础动作降低失败概率的幅度为
In situation $x$, the candidate's reduction in failure probability relative to the base action is
$g_t(x)=Q_t^0(x,B)-Q_t^0(x,\widehat u)$.
% 由于任务后果是失败概率，$g_t(x)\in[-1,1]$。
Because the compared outcomes are failure probabilities, $g_t(x)\in[-1,1]$.
% 有效收益为采用后的平均任务收益，取决于哪些修正及时采用，而非仅总体采用率。
The effective benefit is $\Gamma_{\mathrm{exec},\mu}=\E_\mu[\Lambda g_t(x)]$;
it depends on which corrections are adopted in time, not merely the overall adoption rate.

% 参照保留相同计算与接收流程，以符合共同过程条件；直接关闭反馈计算可能改变基础计划时序，是不同的系统级比较。
The output-masked reference retains the same computation and reception process, as required by Assumption~\ref{ass:commonprocess}.
Skipping feedback computation can change base-plan timing under the shared-worker schedule and therefore defines a different system-level comparison.

% \begin{proposition}[有效反馈与任务完成]
\begin{proposition}[Effective feedback and task completion]
\label{prop:executedtaskgain}
% 在假设~\ref{ass:commonprocess}下，
Under Assumption~\ref{ass:commonprocess},
% 若完整部署策略 $\pi$ 在其访问历史上执行上述 $u_{\mathrm{exec}}$，
if the complete deployed policy $\pi$ executes the above $u_{\mathrm{exec}}$ on its visited histories,
% 而参考策略在相同历史上执行 $B=\pi_0(x)$，则
while the reference policy executes $B=\pi_0(x)$ on the same histories, then
\begin{equation}
 J(\pi)-J(\pi_0)
 =T\,\E_{\bar\mu_\pi}[\Lambda g_t(x)].
 \label{eq:closedlooptaskgain}
\end{equation}
% 因此，实际访问分布上的有效任务收益为正，当且仅当任务完成概率高于参考策略。
Thus, the effective task benefit under the actual visitation distribution is positive if and only if task-completion probability exceeds that of the reference policy.
\end{proposition}

% 进一步，若 $Z_t^{\mathrm{fb}}$ 可读出执行时刻关系 $G_t$ 的估计 $\widehat G_t$ 与上下文 $\chi_t$，
Furthermore, if $Z_t^{\mathrm{fb}}$ admits readouts of an estimate $\widehat G_t$ of the execution-time relations $G_t$ and context $\chi_t$,
% 且这些量在 $\bar\mu_\pi$ 上满足假设~\ref{ass:taskrelations}，则可应用C的关系界。
and these quantities satisfy Assumption~\ref{ass:taskrelations} under $\bar\mu_\pi$, the relation bound in Section~\ref{sec:representationanalysis} applies.
% 此处 $\delta_G=\E_{\bar\mu_\pi}\norm{G_t-\widehat G_t}$ 及其余关系常数均针对反馈表示与执行时刻；
Here, $\delta_G=\E_{\bar\mu_\pi}\norm{G_t-\widehat G_t}$ and all other relation constants concern the feedback representation and execution time.
% 计划时刻 $\tau_p$ 的查询精度不能替代它们（附录~\ref{app:feedbackrelationtime}）。
Query accuracy at the planning time $\tau_p$ cannot replace these quantities (Appendix~\ref{app:feedbackrelationtime}).
% 以 $p_{\mathrm{miss}}=\Prob_{\bar\mu_\pi}(\Lambda=0)$ 表示未采用概率，结合时效界得到
Let $p_{\mathrm{miss}}=\Prob_{\bar\mu_\pi}(\Lambda=0)$ denote the probability of nonadoption. Combining this with the timing bound gives
\begin{equation}
 \begin{aligned}
 J(\pi)-J(\pi_0)\ge T\big[&
 \epsilon_{\mathrm{base}}+\Delta_{\mathrm{new}}+\Delta_{\mathrm{hist}}\\
 &-2(\epsilon_{\mathrm{rel}}+K_G\delta_G)\\
 &-\epsilon_{\mathrm{dec}}(\widehat u)-p_{\mathrm{miss}}\big].
 \end{aligned}
 \label{eq:endtoendtaskbound}
\end{equation}
% 右端为正时，信息价值及基础决策的改进空间超过表示、实现和时效损失，
A positive right-hand side means that information value and the room for improving the base decision exceed representation, realization, and timing losses,
% 从而给出任务完成概率提高的充分条件。
providing a sufficient condition for increased task-completion probability.
% 所有量在同一部署分布评价；这是固定基础策略上的反馈增益，不能将不同策略分布上的局部收益相加。
All terms are evaluated under the same deployment distribution for a fixed base policy.
If $\pi_0$ already includes semantics, this is a feedback-increment guarantee, not the full method's gain over RGB; local gains under different policy distributions cannot simply be added.
% 若关系条件仅在可见等有效子域成立，关系界只适用于该条件分布；
If the relation conditions hold only on a valid subdomain, such as visible cases, the relation bound applies only to that conditional distribution.
% 总体收益还须计入无效子域的贡献，不能直接由式\eqref{eq:endtoendtaskbound}获得保证。
The overall benefit must also account for the invalid subdomain and cannot be guaranteed directly by~\eqref{eq:endtoendtaskbound}.

% 动作与几何监督本身不确保该充分条件成立；回放迁移与动作诊断见附录。
Action and geometric supervision alone do not establish this sufficient condition; replay-transfer bounds and action diagnostics are provided in Appendix~\ref{app:feedbackproofs}.
